%% ============================================================
%% 摘要（中文版）
%% 结构：① 研究背景+问题引出 → ② 本文内容 QCCK → ③ QCCK 基础上提出 QCCK-M
%%       ④ 实验（Beijing 11% + 可解释性） → ⑤ 研究意义 for future
%% 纪律：缩写英文全称只在摘要首次出现时给出一次；
%%       QCCK-M 的英文全称是 QCCK-Model，不是 QCCK-Mamba。
%% 编译：xelatex（ctex）
%% ============================================================

\section*{摘要}

（第一句研究背景）
针对这一问题，本文提出一种量子跨变量耦合核（Quantum Cross-Variable Coupling Kernel，QCCK），利用频域对齐、量子编码与量子核三个环节显式构造变量之间的耦合关系。
在此基础上，本文进一步构建面向多变量时间序列预测的模型 QCCK-Model（QCCK-M），整合输入预处理、主干编码器的逐层更新、跨变量耦合融合与预测输出，把 QCCK 嵌入主干编码器层之间，使显式度量得到的耦合矩阵直接参与跨变量信息融合。
多组实验结果表明，QCCK-M 在多个公开数据集上优于现有方法，其中在 Beijing 数据集上 MSE 降低 11\%；同时，耦合矩阵以显式数值呈现变量间的耦合强度，使模型所依据的变量关系可以被直接读取。
上述结果表明，QCCK 为多变量时间序列预测中的跨变量耦合建模提供了一种可显式度量、可解释的途径，有望为量子核方法在时间序列任务中的应用提供新的思路。
